% !TEX program = xelatex
% Lernplatt - by Can Siebert
% Kurs 22 (Bo 143) - Tag 9 - Lehrskript
% Performance-Marketing im Vertrieb - Anzeigen, Conversion & Landing Pages
%
% Self-contained xelatex source. Kein biblatex, keine externen Bilder.

\documentclass[11pt,a4paper,oneside]{article}

% --- Encoding & Sprache --------------------------------------------------
\usepackage{polyglossia}
\setdefaultlanguage{german}

% --- Geometrie -----------------------------------------------------------
\usepackage[a4paper,top=2cm,bottom=2cm,outer=2.5cm,inner=2.5cm,bindingoffset=1.5cm]{geometry}

% --- Fonts ---------------------------------------------------------------
\usepackage{fontspec}
\defaultfontfeatures{Ligatures=TeX}

\IfFontExistsTF{Charter}{%
  \setmainfont{Charter}[%
    Numbers=OldStyle,%
    UprightFeatures   = {SmallCapsFeatures={Letters=SmallCaps}}%
  ]%
}{%
  \IfFontExistsTF{Bitstream Charter}{%
    \setmainfont{Bitstream Charter}%
  }{%
    \setmainfont{TeX Gyre Schola}%
  }%
}

\IfFontExistsTF{Source Sans 3}{%
  \setsansfont{Source Sans 3}%
}{%
  \IfFontExistsTF{Source Sans Pro}{%
    \setsansfont{Source Sans Pro}%
  }{%
    \setsansfont{TeX Gyre Heros}%
  }%
}

\newcommand{\caveatfontfallback}{\itshape}
\IfFontExistsTF{Caveat}{%
  \newfontfamily\caveatfont{Caveat}[Scale=1.15]%
}{%
  \newcommand\caveatfont{\caveatfontfallback}%
}

\setmonofont{Latin Modern Mono}

% --- Mikro-Typografie ----------------------------------------------------
\usepackage{microtype}
\linespread{1.15}

% --- Farben & Boxen ------------------------------------------------------
\usepackage{xcolor}
\definecolor{lpInk}{RGB}{20,28,38}
\definecolor{lpAccent}{RGB}{22,90,140}
\definecolor{lpAccentLight}{RGB}{210,228,240}
\definecolor{lpAmber}{RGB}{222,138,30}
\definecolor{lpAmberLight}{RGB}{255,243,217}
\definecolor{lpAmberBorder}{RGB}{196,118,18}
\definecolor{lpGray}{RGB}{96,104,114}
\definecolor{lpGrayLight}{RGB}{238,240,243}

\usepackage[most]{tcolorbox}
\tcbuselibrary{breakable,skins}

\newtcolorbox{eselsbox}[1][Eselsbrücke]{%
  enhanced, breakable,
  colback=lpAmberLight,
  colframe=lpAmberBorder,
  boxrule=0.6pt,
  arc=2pt,
  borderline={0.4pt}{0pt}{lpAmberBorder, dashed},
  left=10pt,right=10pt,top=8pt,bottom=8pt,
  fonttitle=\sffamily\bfseries\color{lpAmberBorder},
  coltitle=lpAmberBorder,
  title={#1},
  attach boxed title to top left={xshift=8pt,yshift=-8pt},
  boxed title style={size=small, colback=lpAmberLight, colframe=lpAmberBorder, boxrule=0.4pt, arc=1pt},
  top=14pt
}

\newtcolorbox{merkbox}[1][Merksatz]{%
  enhanced, breakable,
  colback=lpAccentLight,
  colframe=lpAccent,
  boxrule=0.6pt,
  arc=2pt,
  left=10pt,right=10pt,top=8pt,bottom=8pt,
  fonttitle=\sffamily\bfseries\color{white},
  coltitle=white,
  title={#1},
  attach boxed title to top left={xshift=8pt,yshift=-8pt},
  boxed title style={size=small, colback=lpAccent, colframe=lpAccent, boxrule=0.4pt, arc=1pt},
  top=14pt
}

% --- Listen --------------------------------------------------------------
\usepackage{enumitem}
\setlist{itemsep=2pt, topsep=4pt, parsep=0pt}
\setlist[itemize,1]{label={\textbullet},leftmargin=1.6em}
\setlist[itemize,2]{label={\textendash},leftmargin=1.4em}
\setlist[enumerate,1]{label=\arabic*., leftmargin=1.8em}

% --- Tabellen ------------------------------------------------------------
\usepackage{tabularx}
\usepackage{array}
\usepackage{booktabs}
\usepackage{longtable}

% --- Kopf-/Fusszeilen ----------------------------------------------------
\usepackage{fancyhdr}
\usepackage{lastpage}
\pagestyle{fancy}
\fancyhf{}
\renewcommand{\headrulewidth}{0.3pt}
\renewcommand{\footrulewidth}{0pt}
\fancyhead[L]{\footnotesize\sffamily\color{lpGray}Lernplatt \textperiodcentered{} Kurs 22 \textperiodcentered{} Tag 9}
\fancyhead[R]{\footnotesize\sffamily\color{lpGray}Performance-Marketing im Vertrieb}
\fancyfoot[L]{\footnotesize\sffamily\color{lpGray}\textcopyright{} Can Siebert}
\fancyfoot[R]{\footnotesize\sffamily\color{lpGray}\thepage{} / \pageref{LastPage}}

\fancypagestyle{plain}{%
  \fancyhf{}%
  \renewcommand{\headrulewidth}{0pt}%
  \fancyfoot[C]{\footnotesize\sffamily\color{lpGray}\thepage{} / \pageref{LastPage}}%
}

% --- Section-Styling -----------------------------------------------------
\usepackage[explicit]{titlesec}

\titleformat{\section}[block]
  {\sffamily\Large\bfseries\color{lpAccent}}
  {\thesection.}{0.6em}{#1}
  [{\vspace{2pt}\color{lpAccent}\titlerule[0.6pt]}]

\titleformat{\subsection}[block]
  {\sffamily\large\bfseries\color{lpInk}}
  {\thesubsection}{0.6em}{#1}

\titleformat{\subsubsection}[block]
  {\sffamily\normalsize\bfseries\color{lpInk}}
  {}{0pt}{#1}

\titlespacing*{\section}{0pt}{14pt}{8pt}
\titlespacing*{\subsection}{0pt}{10pt}{4pt}
\titlespacing*{\subsubsection}{0pt}{8pt}{2pt}

% --- Hyperlinks ----------------------------------------------------------
\usepackage[hidelinks,unicode]{hyperref}

% --- TOC -----------------------------------------------------------------
\renewcommand{\contentsname}{\sffamily\Large\bfseries\color{lpAccent}Inhalt}

% --- Helfer --------------------------------------------------------------
\newcommand{\akzent}[1]{{\caveatfont\color{lpAmber}#1}}
\newcommand{\fachbegriff}[1]{\textbf{#1}}
\newcommand{\quelle}[1]{{\footnotesize\color{lpGray}(#1)}}

% =========================================================================
\begin{document}

% --- COVER ---------------------------------------------------------------
\begin{titlepage}
  \pagestyle{empty}
  \vspace*{1.5cm}
  \noindent{\sffamily\footnotesize\color{lpGray}LERNPLATT \textperiodcentered{} BY CAN SIEBERT}\\[2pt]
  \noindent{\color{lpAccent}\rule{\linewidth}{1.2pt}}\\[4pt]
  \noindent{\sffamily\footnotesize\color{lpGray}MODUL 5 \textperiodcentered{} TAG 9 \textperiodcentered{} KURS 22 (Bo 143) \textperiodcentered{} 10.\,Juni 2026}

  \vspace{3cm}
  {\sffamily\fontsize{30}{34}\selectfont\bfseries\color{lpInk}Performance-\\Marketing\\im Vertrieb\par}

  \vspace{0.7cm}
  {\sffamily\large\color{lpGray}Anzeigen, Landing Pages \& Conversion --\\messbar planen, qualifiziert wachsen.\par}

  \vspace{1.5cm}
  {\caveatfont\LARGE\color{lpAmber}Lehrskript -- Tag 9 von 16}\par

  \vfill

  \noindent\begin{minipage}{0.55\linewidth}
    {\sffamily\footnotesize\color{lpGray}Lernpfad:}\\
    {\sffamily\small\color{lpInk}Wissen \textrightarrow{} Anwendung \textrightarrow{} Management}\\[10pt]
    {\sffamily\footnotesize\color{lpGray}Bearbeitungsdauer:}\\
    {\sffamily\small\color{lpInk}1 Kurstag}
  \end{minipage}\hfill
  \begin{minipage}{0.4\linewidth}
    \raggedleft
    {\sffamily\footnotesize\color{lpGray}Dozent:}\\
    {\sffamily\small\color{lpInk}Can Siebert}\\[10pt]
    {\sffamily\footnotesize\color{lpGray}Format:}\\
    {\sffamily\small\color{lpInk}Theorie \textperiodcentered{} Fallstudie \textperiodcentered{} Coaching}
  \end{minipage}

  \vspace{0.5cm}
  \noindent{\color{lpAccent}\rule{\linewidth}{0.6pt}}
\end{titlepage}

% --- LERNZIELE -----------------------------------------------------------
\section*{Lernziele dieses Tages}
\addcontentsline{toc}{section}{Lernziele dieses Tages}
\thispagestyle{plain}

\noindent Performance-Marketing wird in vielen Organisationen als ``Werbung mit Tracking'' verstanden. Das verkürzt es. Performance-Marketing ist eine \emph{Vertriebsdisziplin}: datenbasiert, hypothesengetrieben, am Vertriebsergebnis ausgerichtet. Heute lernen Sie, wie Anzeigen, Landing Pages und Conversion-Optimierung als zusammenhängendes System zur Leadgenerierung beitragen -- und welche Entscheidungen auf Senior-Level fallen.

\subsection*{L1 -- Wissen (Reproduktion und Einordnung)}
\begin{itemize}
  \item Grundlagen des Performance-Marketings im Vertrieb verstehen.
  \item Unterschied zwischen Reichweite, Traffic, Leads, Conversion und Umsatz kennen.
  \item Zielgruppen und Buyer Personas für Kampagnen strukturieren.
  \item Grundlagen strategischer Anzeigenplanung verstehen.
  \item Landing Pages als Conversion-Instrument einordnen.
\end{itemize}

\subsection*{L2 -- Anwendung (Transfer auf konkrete Situationen)}
\begin{itemize}
  \item Kampagnenziele aus Vertriebszielen ableiten.
  \item Zielgruppen und Botschaften für Anzeigen entwickeln.
  \item Anzeigenkanäle nach Zielgruppenpassung, Kosten, Datenwirkung und Conversion-Potenzial bewerten.
  \item Landing Pages hinsichtlich Nutzerführung, Vertrauen und Abschlussorientierung analysieren.
  \item Conversion-Hürden erkennen und Optimierungsmaßnahmen ableiten.
\end{itemize}

\subsection*{L3 -- Management (Entscheidung und Verantwortung)}
\begin{itemize}
  \item Performance-Marketing als Investitions- und Steuerungsentscheidung bewerten.
  \item Zielkonflikte zwischen Leadmenge, Leadqualität, Kosten, Geschwindigkeit und Vertriebsressourcen managen.
  \item Entscheidungen unter Unsicherheit, Budgetdruck und Zeitdruck treffen.
  \item Verantwortung für eine skalierbare Performance-Marketing-Architektur übernehmen.
\end{itemize}

\begin{merkbox}[Roter Faden]
Performance-Marketing endet nicht beim Klick. Es endet erst beim \emph{Vertriebsergebnis}: qualifizierter Lead, Demo, Angebot, Abschluss, Deckungsbeitrag. Wer nur bis ``Cost per Click'' misst, betreibt Werbung. Wer bis Deckungsbeitrag und CLV misst, betreibt Vertriebssteuerung.
\end{merkbox}

\clearpage

% --- 1. EINSTIEG ---------------------------------------------------------
\section{Einstieg: Was Performance-Marketing wirklich ist}

\emph{Performance-Marketing} ist eine datenbasierte, ergebnisorientierte Marketingdisziplin, in der jede Maßnahme messbar an einem konkreten Vertriebsziel ausgerichtet ist. Anzeigen, Landing Pages, Tracking, Conversion-Optimierung und Funnel-Analyse bilden ein zusammenhängendes System. \quelle{Chaffey \& Ellis-Chadwick, 2022; Ryan, 2020}

\subsection{Drei Verschiebungen gegenüber klassischem Marketing}

\begin{enumerate}
  \item \fachbegriff{Vom Branding zum Vertriebsergebnis.} Klassisches Marketing baut Marke; Performance-Marketing erzeugt qualifizierte Nachfrage und nachweisbar Pipeline.
  \item \fachbegriff{Vom Streuverlust zur Zielgruppensteuerung.} Anzeigen lassen sich auf Branche, Rolle, Größe, Verhalten zielgenau steuern -- besonders im B2B über LinkedIn.
  \item \fachbegriff{Von Aktivität zu Hypothese.} Eine Kampagne ist eine Hypothese, die getestet wird. Was nicht funktioniert, wird gestoppt; was funktioniert, wird skaliert.
\end{enumerate}

\subsection{Die typische Performance-Falle}

Die häufigste Falle: Performance-Marketing wird isoliert von Vertriebskapazität und Landing Page geplant. ``Wir schalten Anzeigen'' -- und die Leads landen in einem Vertriebsteam, das sie nicht bearbeiten kann, oder auf einer Landing Page, die nicht konvertiert. \quelle{Ellis \& Brown, 2017}

\subsection{Was Tag~9 leistet}

Wir bauen das System auf: Begriffsklärung (Sektion 2), Funnel-Logik (Sektion 3), Zielgruppen und Personas (Sektion 4), Anzeigenkanäle (Sektion 5), Botschaftenlogik (Sektion 6), Landing Pages (Sektion 7), Conversion-Optimierung (Sektion 8) und Fallstudie ``DataSecure Solutions'' (Sektion 9). Sektion 10 behandelt die Management-Perspektive.

\begin{eselsbox}[Eselsbrücke: \akzent{A-L-C-V} -- vier Phasen im Performance-Funnel]
Vier Buchstaben, vier Phasen, ein Funnel:\\[2pt]
\textbf{A}nzeige -- die Aufmerksamkeit ziehen.\\
\textbf{L}anding Page -- den Klick zum Lead wandeln.\\
\textbf{C}onversion / Leadqualifizierung -- Vertrieb übernimmt.\\
\textbf{V}ertriebsergebnis -- Demo, Angebot, Abschluss.\\[2pt]
\emph{Wenn eine der vier Stufen wackelt, verpufft das ganze Budget.}
\end{eselsbox}

% --- 2. BEGRIFFE -- REICHWEITE BIS UMSATZ --------------------------------
\section{Vom Klick zum Umsatz}

Eine zentrale Disziplin im Performance-Marketing ist saubere Begriffsabgrenzung. Wer die Größen nicht trennt, optimiert die falschen Schritte.

\subsection{Sechs zentrale Größen}

\begin{enumerate}
  \item \fachbegriff{Impressionen.} Wie oft wurde unsere Anzeige angezeigt?
  \item \fachbegriff{Klicks.} Wie viele Personen haben darauf reagiert? (Klickrate / CTR = Klicks / Impressionen.)
  \item \fachbegriff{Traffic.} Wie viele Personen sind auf der Landing Page gelandet?
  \item \fachbegriff{Leads.} Wie viele haben das Conversion-Ereignis ausgelöst (Formularabsendung, Demo-Anfrage, Download)?
  \item \fachbegriff{Qualifizierte Leads.} Wie viele der Leads passen zur Zielgruppe und zeigen echte Kaufabsicht? (Marketing-Qualified / Sales-Qualified Leads.)
  \item \fachbegriff{Umsatz / Deckungsbeitrag.} Wie viel Geschäft entsteht aus den qualifizierten Leads?
\end{enumerate}

\subsection{Zentrale KPIs entlang des Funnels}

\begin{itemize}
  \item \fachbegriff{Klickrate (CTR).} Klicks / Impressionen.
  \item \fachbegriff{Kosten pro Klick (CPC).} Werbekosten / Klicks.
  \item \fachbegriff{Conversion Rate.} Leads / Traffic (oder Lead / Klick).
  \item \fachbegriff{Kosten pro Lead (CPL).} Werbekosten / Leads.
  \item \fachbegriff{Kosten pro qualifiziertem Lead.} Werbekosten / qualifizierte Leads.
  \item \fachbegriff{Customer Acquisition Cost (CAC).} Werbe- + Vertriebskosten / gewonnene Kunden.
  \item \fachbegriff{Return on Ad Spend (ROAS).} Umsatz aus Werbung / Werbeausgaben.
  \item \fachbegriff{Customer Lifetime Value (CLV).} Erwarteter Wert eines Kunden über die Beziehung.
\end{itemize}

\subsection{Welche Größe für welches Reporting?}

\begin{tabularx}{\linewidth}{@{}lXX@{}}
\toprule
\textbf{Ebene} & \textbf{Wichtige KPIs} & \textbf{Entscheidungsfokus} \\
\midrule
Operativ (Kampagne) & CTR, CPC, Conversion Rate & Anzeigen, Landing Page optimieren \\
Taktisch (Kanal) & CPL, qualifizierte Leads, Demo-Rate & Budget umverteilen, Kanal stoppen oder skalieren \\
Strategisch (Vertrieb) & CAC, CLV, ROAS, Deckungsbeitrag & Investitions- und Kanal-Mix-Entscheidung \\
\bottomrule
\end{tabularx}

\begin{merkbox}[Funnel-Faustregel]
Optimieren Sie nie isoliert eine Funnel-Stufe. Eine erhöhte Klickrate ist wertlos, wenn die Conversion sinkt. Eine niedrige Conversion ist nicht zwingend Schuld der Landing Page -- sie kann an falsch ausgesteuerten Anzeigen liegen. Erst das gemeinsame Bild macht Sinn.
\end{merkbox}

% --- 3. FUNNEL-LOGIK -----------------------------------------------------
\section{Funnel-Logik im Vertrieb}

Ein \fachbegriff{Funnel} ist die strukturierte Abfolge der Schritte vom ersten Kontakt bis zum Abschluss. Im B2B-Performance-Marketing umfasst er typischerweise acht Stufen. \quelle{Chaffey \& Ellis-Chadwick, 2022; Ellis \& Brown, 2017}

\subsection{Acht Stufen im B2B-Performance-Funnel}

\begin{enumerate}
  \item \textbf{Anzeige} -- Impressionen, Targeting, Botschaft.
  \item \textbf{Klick} -- Person reagiert auf Anzeige.
  \item \textbf{Landing Page} -- gezielte Seite, die das Anzeigeversprechen einlöst.
  \item \textbf{Conversion} -- Formularabsendung, Download, Demo-Anfrage.
  \item \textbf{Leadqualifizierung} -- Marketing prüft Score / Vertrieb prüft Eignung.
  \item \textbf{Vertriebskontakt} -- Erstgespräch, Bedarfsanalyse.
  \item \textbf{Angebot} -- konkretes Angebot, oft individuell.
  \item \textbf{Abschluss} -- Vertrag, Kunde gewonnen.
\end{enumerate}

\subsection{Wo entstehen die größten Verluste?}

Empirisch ist gut belegt, dass die größten Verluste meist an den \emph{Übergängen} entstehen: Klick \textrightarrow{} Conversion (schlechte Landing Page) und Lead \textrightarrow{} Vertriebskontakt (zu langsame Reaktion, kein Lead Scoring). Wer hier 20--30\,\% verbessert, vervielfacht effektiv das Marketingbudget. \quelle{Ash et al., 2012}

\subsection{Funnel als gemeinsame Sicht}

Ein zentrales Senior-Prinzip: Marketing, Vertrieb und Controlling sollten den \emph{gleichen Funnel} sehen. Wenn Marketing CTR und CPL berichtet, Vertrieb Abschlüsse und Pipeline -- und beide nicht miteinander reden -- bleibt jede Optimierung Stückwerk. Erst ein gemeinsames Reporting (mit Lead-, Demo-, Angebots-, Abschlussrate) erzeugt Steuerungsfähigkeit.

\begin{eselsbox}[Eselsbrücke: \akzent{A-K-L-C-Q-V-A-A}]
Acht Stufen, die jede B2B-Performance-Kampagne durchläuft:\\[2pt]
\textbf{A}nzeige -- \textbf{K}lick -- \textbf{L}anding Page -- \textbf{C}onversion -- \textbf{Q}ualifizierung -- \textbf{V}ertriebskontakt -- \textbf{A}ngebot -- \textbf{A}bschluss.\\[2pt]
\emph{Wer als Marketing-Verantwortlicher nur bis ``Q'' denkt und der Vertrieb erst ab ``V'' anfängt, hat in der Mitte den Funnel auf zwei Hälften zerschnitten -- die Leads fallen heraus.}
\end{eselsbox}

% --- 4. ZIELGRUPPEN UND BUYER PERSONAS ---------------------------------
\section{Zielgruppen und Buyer Personas für Kampagnen}

Eine Performance-Kampagne ist nur so gut wie ihre Zielgruppensteuerung. Wer ``allen'' zeigen will, zahlt für irrelevante Klicks und verbrennt Vertriebskapazität. \quelle{Kotler et al., 2017; Homburg, 2020}

\subsection{Sechs Dimensionen der Zielgruppenanalyse}

\begin{enumerate}
  \item \fachbegriff{Bedarf.} Welches konkrete Problem hat die Zielgruppe?
  \item \fachbegriff{Rolle.} Welche Position im Unternehmen (IT-Leiter, Geschäftsführung, Datenschutzbeauftragter)?
  \item \fachbegriff{Kaufmotiv.} Welche Motive treiben (Risikoreduktion, Compliance, Wirtschaftlichkeit, Geschwindigkeit)?
  \item \fachbegriff{Informationsstand.} Ist die Person Anfänger im Thema oder erfahrener Entscheider?
  \item \fachbegriff{Entscheidungsbefugnis.} Entscheidet sie alleine, mit anderen, gar nicht?
  \item \fachbegriff{Schmerzpunkte.} Was nervt sie konkret im Alltag -- und welches Versprechen würde sie hören wollen?
\end{enumerate}

\subsection{Buying Center im Performance-Marketing}

Im B2B entscheiden mehrere Rollen gemeinsam. Eine reife Kampagnenarchitektur adressiert mehrere Rollen -- mit unterschiedlichen Botschaften, aber konsistentem Versprechen.

Für ein Cybersecurity-Angebot:

\begin{itemize}
  \item \textbf{IT-Leiter.} Fachliche Sicherheit, technische Passung, Integrationsfähigkeit.
  \item \textbf{Geschäftsführer.} Risiko, Kosten, Haftung, Geschäftskontinuität.
  \item \textbf{Datenschutzbeauftragter.} Nachweise, Prozesse, Auditfähigkeit.
  \item \textbf{Compliance-Verantwortlicher.} Regulatorische Standards, Dokumentation.
\end{itemize}

\subsection{Buyer Personas konkret}

Eine Buyer Persona verdichtet eine Zielgruppe in einen archetypischen ``Kunden im Kopf'':

\begin{itemize}
  \item Name, Rolle, Branche, Unternehmensgröße.
  \item Ziele und Erfolgsmaßstäbe in der Rolle.
  \item Typische Tagesabläufe und Belastungen.
  \item Informationsquellen (Fachpresse, LinkedIn, Konferenzen).
  \item Häufige Einwände gegen Angebote in unserer Kategorie.
  \item Sprachlich passende Begriffe (technisch vs.\ management-orientiert).
\end{itemize}

Drei bis fünf Personas reichen meist -- mehr sind im Mittelstand selten ehrlich pflegbar.

% --- 5. ANZEIGENKANAELE -------------------------------------------------
\section{Anzeigenkanäle und ihre Logik}

Sechs Anzeigenkanäle dominieren das B2B-Performance-Marketing 2026. Jeder hat eine eigene Logik, eigene Stärken und eigene Fallen.

\subsection{Sechs Kanäle im Vergleich}

\begin{enumerate}
  \item \fachbegriff{Google Ads (Search).} Reagiert auf konkrete Suchanfragen -- ideal bei kaufnaher Intention. Hoher Wettbewerb auf typischen B2B-Begriffen, daher hoher CPC. Sehr gut für ``akute Nachfrage''.
  \item \fachbegriff{LinkedIn Ads.} Targeting nach Branche, Größe, Rolle, Seniorität. Teurer CPC, aber häufig die höchste Leadqualität im B2B. Stark bei ``erklärungsbedürftig'' und ``Buying Center''.
  \item \fachbegriff{Meta Ads (Facebook / Instagram).} B2C-stark, B2B-mäßig -- ausgenommen Produkte mit konsumnaher Anwendung (z.\,B.\ kleine Mittelständler, Selbstständige).
  \item \fachbegriff{Plattformanzeigen.} Anzeigen auf Marktplätzen / Branchenportalen -- nah am Kaufkontext, oft niedrige Streuverluste.
  \item \fachbegriff{Retargeting.} Anzeigen für Personen, die schon auf der Website waren. Hohe Conversion, niedrige Kosten -- aber begrenzte Zielgruppe. \quelle{Lambrecht \& Tucker, 2013}
  \item \fachbegriff{Display-Kampagnen.} Breite Reichweite mit Banner-Anzeigen. Hoher Streuverlust; für B2B-Mittelstand selten effizient.
\end{enumerate}

\subsection{Kanalbewertung im B2B}

\begin{tabularx}{\linewidth}{@{}lXXXX@{}}
\toprule
\textbf{Kanal} & \textbf{Zielgruppenpassung} & \textbf{Kosten} & \textbf{Leadqualität} & \textbf{Geschwindigkeit} \\
\midrule
Google Ads (Search) & hoch (bei aktiver Suche) & mittel--hoch & mittel--hoch & sehr schnell \\
LinkedIn Ads & sehr hoch (Rollen-Targeting) & hoch & hoch & mittel \\
Plattformanzeigen & hoch (im Branchenkontext) & mittel & hoch & schnell \\
Retargeting & sehr hoch (Vorinteresse) & niedrig & hoch & sehr schnell \\
Display & niedrig (Streuung) & niedrig pro Klick & niedrig & schnell \\
Meta Ads & B2C-orientiert & niedrig--mittel & niedrig--mittel & schnell \\
\bottomrule
\end{tabularx}

\subsection{Wann ein Kanal \emph{nicht} sinnvoll ist}

\begin{itemize}
  \item Display-Kampagnen bei sicherheits-/komplexitätsbedürftigen B2B-Angeboten -- Reichweite ohne Vertriebskapazität gebunden zu sehen, ist teure Sichtbarkeit.
  \item Meta Ads bei reinen B2B-Enterprise-Produkten.
  \item LinkedIn Ads bei klassischen B2C-Konsumgütern (zu teuer für die Conversion-Ökonomie).
\end{itemize}

\begin{merkbox}[Kanal-Faustregel]
Wählen Sie Anzeigenkanäle aus zwei Achsen: \emph{wo} ist die Zielgruppe -- und \emph{in welcher Phase} der Customer Journey ist sie dort? Google Ads = akute Suche. LinkedIn Ads = Buying Center erreichen. Retargeting = Vorinteresse vertiefen. Wer eine Phase ohne passendem Kanal bespielt, verbrennt Budget.
\end{merkbox}

% --- 6. BOTSCHAFTEN -----------------------------------------------------
\section{Botschaften entwickeln: Problem, Nutzen, Beweis, CTA}

Eine wirksame Anzeige beantwortet vier Fragen -- in der Reihenfolge, in der sie der Leser stellt: \quelle{Ryan, 2020}

\subsection{Vier Bausteine einer Botschaft}

\begin{enumerate}
  \item \fachbegriff{Problem.} Erkenne ich mich in dieser Beschreibung wieder? (``Cyberrisiken im Mittelstand reduzieren'')
  \item \fachbegriff{Nutzenversprechen.} Was bekomme ich, wenn ich klicke? (``Sicherheitscheck in 30~Minuten, ohne IT-Aufwand'')
  \item \fachbegriff{Beweis.} Warum sollte ich es glauben? (``150+ Mittelständler haben unseren Sicherheitscheck genutzt'' oder Zertifikate)
  \item \fachbegriff{Call-to-Action.} Was soll ich konkret tun? (``Sicherheitscheck anfragen'')
\end{enumerate}

\subsection{Botschaftendifferenzierung pro Rolle}

Pro Buyer Persona eine eigene Botschaft -- aber konsistent in Markenversprechen und Look. Beispiel Cybersecurity:

\begin{itemize}
  \item \textbf{IT-Leiter:} ``Lückenlos Schutz integrieren -- ohne Bestandssysteme zu sprengen.''
  \item \textbf{Geschäftsführer:} ``Cyberrisiken planbar reduzieren -- mit klarer Investitionsentscheidung.''
  \item \textbf{Datenschutzbeauftragter:} ``Auditfähigkeit nachweisbar sichern -- alle Prozesse dokumentiert.''
\end{itemize}

\subsection{Botschaft \emph{passt zu} Landing Page}

Eine zentrale Disziplin: die Anzeige muss zur Landing Page \emph{passen}. Wer mit ``Sicherheitscheck in 30 Minuten'' wirbt, und der Klick führt auf eine generische Homepage mit Firmengeschichte, verliert den Lead. Das passt-Prinzip ist kein Detail, sondern der wichtigste einzelne Conversion-Hebel.

% --- 7. LANDING PAGES ---------------------------------------------------
\section{Landing Pages: das Conversion-Instrument}

Eine \fachbegriff{Landing Page} ist eine speziell für eine Kampagne gestaltete Seite -- nicht die Homepage. Sie hat eine Aufgabe: die Conversion zu erreichen. Empirisch ist sie der wichtigste Hebel zwischen Klick und Lead. \quelle{Ash et al., 2012; Ryan, 2020}

\subsection{Acht Elemente einer wirksamen Landing Page}

\begin{enumerate}
  \item \fachbegriff{Klare Überschrift.} Mit Problembezug, in maximal 8--10 Wörtern. (``Cyberrisiken im Mittelstand erkennen und gezielt reduzieren.'')
  \item \fachbegriff{Kurzes Nutzenversprechen.} 1--2 Sätze, was die Person konkret bekommt.
  \item \fachbegriff{Zielgruppenspezifische Abschnitte.} Pro Rolle ein eigener Absatz oder Block.
  \item \fachbegriff{Vertrauenselemente.} Referenzen, Kundenstimmen, Zertifizierungen, Sicherheitsnachweise.
  \item \fachbegriff{Kurze Prozessdarstellung.} Wie läuft der nächste Schritt? (``1. Anfrage senden. 2.\ Kostenloser 30-Min-Check. 3.\ Empfehlungsbericht.'')
  \item \fachbegriff{Call-to-Action.} Klarer, sichtbar, mehrfach platziert. Wortwahl konkret (``Sicherheitscheck anfragen'', nicht ``Absenden'').
  \item \fachbegriff{Formular.} Mit wenigen Pflichtfeldern (3--5), optionale Felder erlaubt.
  \item \fachbegriff{FAQ.} Antworten auf typische Einwände (Aufwand, Kosten, Datenschutz, Ablauf).
\end{enumerate}

\subsection{Conversion-Hürden -- die typischen Bremser}

\begin{itemize}
  \item \textbf{Unklare Botschaft.} Was bekomme ich hier, in einem Satz? Wenn der Leser es nach 5 Sekunden nicht weiß, springt er ab.
  \item \textbf{Zu viele Informationen.} Wer alles erklärt, lenkt vom Conversion-Ziel ab.
  \item \textbf{Fehlendes Vertrauen.} Keine Referenzen, keine Logos, keine Zertifikate -- besonders im B2B kritisch.
  \item \textbf{Lange Formulare.} Jedes zusätzliche Pflichtfeld senkt die Conversion messbar.
  \item \textbf{Widersprüche zwischen Anzeige und Page.} Anzeige verspricht X, Page redet von Y.
  \item \textbf{Langsame Ladezeit.} Bereits ab 3--4 Sekunden Verzögerung steigt die Absprungrate signifikant.
  \item \textbf{Fehlende Mobile-Tauglichkeit.} Im B2B oft unterschätzt -- ein Großteil der Klicks kommt mobil.
  \item \textbf{Fehlende Relevanz zur Anzeige.} ``Bait \& Switch'' bestraft die Plattform mit höherem CPC.
\end{itemize}

\subsection{A/B-Tests und iterative Optimierung}

Landing Pages werden nicht ``richtig'' gebaut -- sie werden \emph{verbessert}. Ein \fachbegriff{A/B-Test} testet zwei Varianten parallel (z.\,B.\ Überschrift A vs.\ Überschrift B) und entscheidet datenbasiert. Eine reife Performance-Organisation hat einen kontinuierlichen Test-Rhythmus (z.\,B.\ ein Test pro Landing Page pro 4 Wochen).

\begin{merkbox}[Landing-Page-Kernsatz]
Die Landing Page ist kein Nebenprodukt der Werbung. Sie ist der wichtigste einzelne Conversion-Hebel. Wer 50.000\,EUR in Anzeigen investiert und 500\,EUR in die Landing Page, hat die Verhältnisse falsch gesetzt.
\end{merkbox}

% --- 8. CONVERSION-OPTIMIERUNG ------------------------------------------
\section{Conversion-Optimierung als laufende Disziplin}

\fachbegriff{Conversion-Rate-Optimization (CRO)} ist die systematische Verbesserung der Conversion entlang des Funnels. Sie unterscheidet sich von ``Marketing'' insofern, dass sie hypothesen- und testgetrieben arbeitet. \quelle{Ash et al., 2012; Ellis \& Brown, 2017}

\subsection{Vier Bausteine einer CRO-Praxis}

\begin{enumerate}
  \item \textbf{Hypothese.} ``Wenn wir das Formular von 12 auf 5 Felder kürzen, steigt die Conversion um mindestens 30\,\%.''
  \item \textbf{Test.} Variante A (aktuell) vs.\ Variante B (Hypothese) parallel über ausreichend Volumen.
  \item \textbf{Auswertung.} Statistische Signifikanz prüfen -- nicht nach 50 Klicks entscheiden.
  \item \textbf{Roll-out oder Verwerfen.} Bei Gewinn übernehmen, bei Verlust verwerfen -- in jedem Fall das gelernte dokumentieren.
\end{enumerate}

\subsection{Was zuerst optimieren?}

Eine pragmatische Reihenfolge:

\begin{enumerate}
  \item \textbf{Botschaftspassung.} Anzeige und Landing Page stimmig.
  \item \textbf{Vertrauen.} Referenzen, Zertifikate, Kundenstimmen sichtbar.
  \item \textbf{Formular.} Pflichtfelder reduzieren, Optionale erlauben.
  \item \textbf{Call-to-Action.} Wortwahl, Platzierung, Wiederholung.
  \item \textbf{Vertrauenshemmer abbauen.} FAQ zu Datenschutz, Kosten, Aufwand.
  \item \textbf{Geschwindigkeit / Mobile.} Technische Hürden beseitigen.
\end{enumerate}

\subsection{Conversion vs.\ Leadqualität}

Eine erhöhte Conversion-Rate ist nicht automatisch gut. Wer das Formular auf zwei Felder kürzt, bekommt mehr Leads -- aber möglicherweise weniger qualifizierte. Daher ist die Steuerungsgröße meistens nicht ``Conversion'', sondern ``Conversion zu qualifiziertem Lead''.

\begin{eselsbox}[Eselsbrücke: \akzent{H-T-A-R}]
Vier Buchstaben, vier Schritte einer Conversion-Optimierung:\\[2pt]
\textbf{H}ypothese -- begründete Vermutung.\\
\textbf{T}est -- A/B im echten Traffic.\\
\textbf{A}uswertung -- mit statistischer Signifikanz.\\
\textbf{R}oll-out oder Verwerfen -- klare Entscheidung.\\[2pt]
\emph{Wer ohne Hypothese testet, sammelt Daten -- nicht Erkenntnisse.}
\end{eselsbox}

% --- 9. FALLSTUDIE-LEITFADEN ---------------------------------------------
\section{Fallstudie-Leitfaden: DataSecure Solutions}

In den Unterrichtseinheiten 4--5 bearbeiten Sie die Fallstudie ``DataSecure Solutions'' -- ein mittelständischer Anbieter von Cybersecurity-Lösungen für Unternehmen mit 100--1.000 Mitarbeitenden, 16\,Mio.\,EUR Umsatz, 8 Vertriebs- und 2 Marketing-Mitarbeitenden.

\subsection{Analytischer Aufbau}

\begin{enumerate}
  \item \textbf{Zielgruppenanalyse.} IT-Leiter, Geschäftsführer, Datenschutzbeauftragte, Compliance -- welche Informationsbedürfnisse?
  \item \textbf{Kampagnenstrategie für 6 Monate.} Welche Kanäle, welche Reihenfolge, welche Ziele?
  \item \textbf{Kanalpriorisierung.} Google, LinkedIn, Webinar, Retargeting, Display -- welche kombiniert man, welche nicht?
  \item \textbf{Landing Page Struktur.} Acht Elemente konkret ausgestalten für ein Demo- oder Sicherheitscheck-Angebot.
  \item \textbf{Conversion-Ziele und KPIs.} Entlang des Funnels.
  \item \textbf{Budgetverteilung.} 150.000\,EUR aufteilen auf Kanäle, Landing Page, Content, Tests.
  \item \textbf{Entscheidung unter Unsicherheit.} Warum bestimmte Kanäle (z.\,B.\ Display) \emph{nicht}?
\end{enumerate}

\subsection{Erwartete Empfehlungsrichtung}

In einer typischen Musterlösung wird priorisiert: LinkedIn Ads (Buying Center), Webinare (Vertrauensaufbau, Fachvermittlung), Google Ads (akute Suchen), Retargeting (für Vor-Interessierte). \emph{Bewusster Verzicht} auf breite Display-Kampagnen, weil sie hohen Budgetverbrauch ohne qualifizierte Leadqualität bedeuten. Die Entscheidung unter Unsicherheit: ``Trotz höherer Reichweite verzichten wir auf breite Display-Kampagnen, weil das Vertriebsteam nur 20 neue Leads pro Woche verarbeiten kann -- die Kombination aus LinkedIn, Google Ads und Webinaren liefert höhere Leadqualität und passt zur Vertriebskapazität.''

\begin{merkbox}[Argumentationsmuster für die Fallstudie]
Eine gute Performance-Marketing-Architektur enthält:
\begin{enumerate}
  \item \textbf{Zielgruppen + Buying Center} -- pro Rolle eine Botschaft.
  \item \textbf{Kanalpriorisierung mit Aufgabe pro Kanal} -- nicht alle Kanäle gleichzeitig.
  \item \textbf{Landing Page Struktur} -- konkret, mit allen acht Elementen.
  \item \textbf{KPI-Architektur entlang des Funnels} -- nicht nur CPC.
  \item \textbf{Bewusster Verzicht auf einen reichweitenstarken Kanal} -- mit Begründung.
  \item \textbf{Annahme + Frühindikator} -- welche Hypothese, welche KPI signalisiert Korrektur?
\end{enumerate}
\end{merkbox}

% --- 10. MANAGEMENT-PERSPEKTIVE ------------------------------------------
\section{Management-Perspektive: Zielkonflikte}

Auf Wissens- und Anwendungsebene sieht Performance-Marketing wie eine sauber strukturierte Kampagnenplanung aus. Auf Managementebene geht es um Zielkonflikte, die nicht ``gelöst'', sondern \emph{entschieden} werden müssen. \quelle{vgl.\ Homburg, 2020; Kaplan \& Norton, 1996}

\subsection{Fünf zentrale Zielkonflikte}

\begin{enumerate}
  \item \textbf{Leadmenge vs.\ Leadqualität.} Mehr Anzeigen erzeugen mehr Leads -- aber häufig schlechtere. Wer den Vertrieb mit unqualifizierten Kontakten überflutet, verbrennt Kapazität.
  \item \textbf{Kosten vs.\ Geschwindigkeit.} Performance-Kanäle wie LinkedIn sind teuer, aber liefern Buying-Center-Reichweite. SEO/Content ist günstig, aber langsam.
  \item \textbf{Reichweite vs.\ Vertriebskapazität.} Mehr Reichweite ist erst dann sinnvoll, wenn der Vertrieb die Leads sinnvoll bearbeiten kann.
  \item \textbf{Kurzfristige Klicks vs.\ Pipeline-Aufbau.} Kampagnen wirken kurzfristig; Content, Vertrauen und Wiederkauf wirken über Quartale.
  \item \textbf{Automatisierung vs.\ Verantwortung.} Algorithmische Anzeigensteuerung ist effizient -- aber sie kann Targeting und Botschaft unsensibel führen, wenn keine Governance dahintersteht.
\end{enumerate}

\subsection{Performance-Marketing als Investitionsentscheidung}

Performance-Marketing ist \emph{keine} Marketingmaßnahme im klassischen Sinn -- es ist eine Investitionsentscheidung über Ressourceneinsatz, Marktprioritäten und Vertriebskapazität. Sie wird auf Senior-Level entschieden, typischerweise durch Chief Revenue Officer oder Geschäftsführung gemeinsam mit Marketingleitung. Wer das Budget allein an die Marketingagentur delegiert, delegiert eine strategische Vertriebsentscheidung.

\subsection{Senior-Level-Test}

Eine senior-level-fähige Performance-Strategie lässt sich in einem Satz erklären, der eine bewusste Verzichtsentscheidung und einen Frühindikator enthält. Beispiel: ``Wir investieren 45\,k in LinkedIn Ads, 30\,k in Google Ads, 25\,k in Webinare und 20\,k in Landing Page \& Tracking -- wir verzichten bewusst auf breite Display-Kampagnen, weil die Vertriebskapazität bei 20 Leads pro Woche begrenzt ist. Wenn die Kosten pro qualifiziertem Lead nach 8 Wochen über 800\,EUR liegen, prüfen wir den Kanal-Mix neu.''

\begin{eselsbox}[Eselsbrücke: \akzent{L-K-R-K-A}]
Fünf Zielkonflikte, die im Performance-Marketing auf Senior-Level entschieden werden:\\[2pt]
\textbf{L}eadmenge vs.\ Leadqualität.\\
\textbf{K}osten vs.\ Geschwindigkeit.\\
\textbf{R}eichweite vs.\ Vertriebskapazität.\\
\textbf{K}urzfristige Klicks vs.\ Pipeline-Aufbau.\\
\textbf{A}utomatisierung vs.\ Verantwortung.\\[2pt]
\emph{Fünf Konflikte, die kein Marketingdienstleister allein verantworten sollte.}
\end{eselsbox}

\begin{merkbox}[Schluss-Merksatz für Tag 9]
Performance-Marketing endet nicht bei Klick-Kosten. Es endet beim Vertriebsergebnis. Die Kampagne ist nur dann erfolgreich, wenn Anzeige, Landing Page, Leadqualifizierung und Vertrieb sauber verbunden sind -- und wenn das Marketing nicht den Vertrieb überfordert, sondern ihn präzise füttert. \emph{Wer Performance-Marketing als Werbung versteht, kauft Klicks. Wer es als Vertriebssteuerung versteht, kauft Kunden.}
\end{merkbox}

% --- 11. LITERATUR -------------------------------------------------------
\clearpage
\section{Literatur}

Im Skript verwendete und für die Vertiefung empfohlene Quellen. Zitierstil: APA-7.

\begin{description}[style=nextline,leftmargin=18pt,labelindent=0pt,itemsep=4pt]
  \item[Ash, T., Ginty, M., \& Page, R. (2012).] \emph{Landing page optimization: The definitive guide to testing and tuning for conversions} (2.\ Aufl.). Sybex.
  \item[Chaffey, D., \& Ellis-Chadwick, F. (2022).] \emph{Digital marketing: Strategy, implementation and practice} (8.\ Aufl.). Pearson.
  \item[Ellis, S., \& Brown, M. (2017).] \emph{Hacking growth: How today's fastest-growing companies drive breakout success}. Crown Business.
  \item[Homburg, C. (2020).] \emph{Marketingmanagement: Strategie -- Instrumente -- Umsetzung -- Unternehmensführung} (7.\ Aufl.). Springer Gabler.
  \item[Kaplan, R.\,S., \& Norton, D.\,P. (1996).] \emph{The balanced scorecard: Translating strategy into action}. Harvard Business School Press.
  \item[Kotler, P., Kartajaya, H., \& Setiawan, I. (2017).] \emph{Marketing 4.0: Moving from traditional to digital}. Wiley.
  \item[Lambrecht, A., \& Tucker, C. (2013).] When does retargeting work? Information specificity in online advertising. \emph{Journal of Marketing Research, 50}(5), 561--576. \href{https://doi.org/10.1509/jmr.11.0503}{https://doi.org/10.1509/jmr.11.0503}
  \item[Ryan, D. (2020).] \emph{Understanding digital marketing: A complete guide to engaging customers with digital marketing} (5.\ Aufl.). Kogan Page.
\end{description}

% --- 12. GLOSSAR ---------------------------------------------------------
\clearpage
\section{Glossar}

Die wichtigsten Begriffe des Tages -- kurz definiert, in alphabetischer Reihenfolge.

\begin{description}[style=nextline,leftmargin=0pt,labelindent=0pt,itemsep=4pt]
  \item[A/B-Test] Vergleich zweier Varianten (Anzeige, Landing Page, Element) im echten Traffic zur datenbasierten Entscheidung.
  \item[Buyer Persona] Archetypische Beschreibung einer Käuferrolle mit Zielen, Schmerzpunkten und Informationsquellen.
  \item[Buying Center] Mehrstufige B2B-Entscheidergruppe (z.\,B.\ IT-Leiter, Geschäftsführer, Compliance) -- jeder benötigt eine eigene Botschaft.
  \item[Call-to-Action (CTA)] Handlungsaufforderung in Anzeige oder Landing Page (z.\,B.\ ``Sicherheitscheck anfragen'').
  \item[Conversion] Messbares Ereignis, das ein Besucher auf der Landing Page auslöst -- Formularabsendung, Download, Demo-Anfrage.
  \item[Conversion Rate] Anteil Besucher, die eine Conversion auslösen.
  \item[Conversion-Hürde] Element, das die Conversion bremst (unklare Botschaft, langes Formular, fehlendes Vertrauen, langsame Ladezeit).
  \item[Conversion-Rate-Optimization (CRO)] Hypothesen- und testgetriebene Verbesserung der Conversion entlang des Funnels.
  \item[Cost per Click (CPC)] Werbekosten geteilt durch Klicks; KPI auf Anzeigenebene.
  \item[Cost per Lead (CPL)] Werbekosten geteilt durch generierte Leads; KPI auf Kanalebene.
  \item[Customer Acquisition Cost (CAC)] Gesamtkosten (Werbung + Vertrieb), die ein gewonnener Kunde verursacht.
  \item[Customer Lifetime Value (CLV)] Erwarteter wirtschaftlicher Wert eines Kunden über die gesamte Beziehung.
  \item[Display-Kampagne] Banner-Anzeigen mit breiter Reichweite -- hoher Streuverlust im B2B.
  \item[Funnel] Geordnete Abfolge der Schritte von der Anzeige bis zum Abschluss; im B2B typischerweise acht Stufen.
  \item[Google Ads] Suchanzeigen, die auf konkrete Suchanfragen ausgeliefert werden; stark bei kaufnaher Intention.
  \item[Impressionen] Anzahl Anzeigen-Einblendungen, ohne Aussage über Wirkung.
  \item[Klickrate (CTR)] Klicks geteilt durch Impressionen; Indikator für Anzeigenrelevanz.
  \item[Landing Page] Speziell für eine Kampagne gestaltete Zielseite -- nicht die Homepage.
  \item[Lead] Interessensbekundung, die in einen Vertriebsprozess überführt werden kann.
  \item[Lead Scoring] Bewertung der Leadqualität anhand definierter Kriterien (Demografie, Verhalten, Bedarfssignale).
  \item[LinkedIn Ads] Anzeigen mit Targeting nach Branche, Rolle, Seniorität; im B2B häufig die höchste Leadqualität.
  \item[Nutzenversprechen] Konkrete Aussage, was der Klick / die Conversion dem Besucher bringt.
  \item[Performance-Marketing] Datenbasierte, ergebnisorientierte Marketingdisziplin entlang eines messbaren Funnels.
  \item[Qualifizierter Lead] Lead, der zur Zielgruppe passt und Kaufabsicht zeigt -- MQL/SQL.
  \item[Retargeting] Anzeigen für Personen, die bereits auf der Website waren; hohe Conversion, niedrige Kosten.
  \item[Return on Ad Spend (ROAS)] Umsatz aus Werbung geteilt durch Werbeausgaben.
  \item[Suchintention] Absicht hinter einem Suchbegriff -- informieren, vergleichen, kaufen, wiederbestellen.
  \item[Vanity-Metrik] Kennzahl ohne Steuerungswirkung (z.\,B.\ reine Impressionen, Followerzahlen).
  \item[Zielkonflikt (Performance)] Spannungsfeld zwischen Leadmenge, Leadqualität, Kosten, Geschwindigkeit, Vertriebskapazität und Pipeline-Aufbau -- bewusst entscheidbar.
\end{description}

\vspace{1cm}
\noindent{\color{lpAccent}\rule{\linewidth}{0.6pt}}\\[4pt]
{\footnotesize\sffamily\color{lpGray}Lernplatt \textperiodcentered{} by Can Siebert \textperiodcentered{} Kurs 22 (Bo 143) \textperiodcentered{} Tag 9 \textperiodcentered{} \textcopyright{} 2026}

\end{document}
